\documentclass[bachelor, och, pract3]{SCWorks}
% параметр - тип обучения - одно из значений:
%    spec     - специальность
%    bachelor - бакалавриат (по умолчанию)
%    master   - магистратура
% параметр - форма обучения - одно из значений:
%    och   - очное (по умолчанию)
%    zaoch - заочное
% параметр - тип работы - одно из значений:
%    referat      - реферат
%    coursework   - курсовая работа (по умолчанию)
%    diploma      - дипломная работа
%    pract2       - отчет по практике, 2 курс, бакалавриат, магистратура
%    pract3       - отчет по практике, 3 курс, бакалавриат
%    pract4       - отчет по практике, 4 курс, бакалавриат
%    practpreddip - отчет по преддипломной практике, 4 курс (= pract3 или nir) 
%    nir          - отчет о научно-исследовательской работе
%    autoref      - автореферат выпускной работы
%    assignment   - задание на выпускную квалификационную работу
%    review       - отзыв руководителя
%    critique     - рецензия на выпускную работу
% параметр - включение шрифта
%    times    - включение шрифта Times New Roman (если установлен)
%               по умолчанию выключен
\usepackage{cmap}
\usepackage[T2A]{fontenc}
\usepackage[utf8]{inputenc}
\usepackage{graphicx}
\usepackage[sort,compress]{cite}
\usepackage{amsmath}
\usepackage{amssymb}
\usepackage{amsthm}
\usepackage{fancyvrb}
\usepackage{longtable}
\usepackage{array}
\usepackage[english,russian]{babel}

% выделение ссылок цветом. чтобы включить- true
\usepackage[colorlinks=false]{hyperref}

\renewcommand{\rmdefault}{cmr} % курсив и полужирный включаются здесь.
\newcommand{\eqdef}{\stackrel {\rm def}{=}}

\newtheorem{lem}{Лемма}

\begin{document}

% Кафедра (в родительном падеже)
\chair{теории функций и стохастического анализа} 
\chairr{кафедра теории функций и стохастического анализа} % для практик!!!!
%\chairr{завод "Тантал"}% для практик!!!!

%Тема работы

\title{Моделирование случайных величин. Вариант 7}

% Курс
\course{2}

% Группа
\group{212}

\department{механико-математического факультета}

% Специальность/направление код - наименование
\napravlenie{01.03.02 "--- Прикладная математика и информатика}
%\napravlenie{38.03.05 "--- Бизнес-информатика}
%\napravlenie{09.04.03 "--- Прикладная информатика}

% Для СТУДЕНТКИ!!!. Для работы студента следующая команда не нужна.
%\studenttitle{Студентки}

% Фамилия, имя, отчество в родительном падеже
\author{Чесакова Максима Евгеньевича}

% Заведующий кафедрой
%\chtitle{д.\,ф.-м.\,н., доцент} % степень, звание
%\chname{С.\,П.\,Сидоров}

%Научный руководитель (для реферата преподаватель проверяющий работу)
\satitle{доцент, к.\,ф.-м.\,н.} %должность, степень, звание
\saname{А.\,Д.\,Луньков} % для 451 гр
%\saname{А.\,К.\,Смирнов} % для 412 гр

% Руководитель практики от организации (только для практики,
% для остальных типов работ не используется)

\patitle{руководитель направления АРМ}
\paname{Д.\,Э.\,Кнутов}

% Семестр (только для практики, для остальных
% типов работ не используется)
\term{5}

% Наименование практики (только для практики, для остальных
% типов работ не используется)
%\practtype{производственная(базовая)}

% Продолжительность практики (количество недель) (только для практики,
% для остальных типов работ не используется)
%\duration{4}

% Даты начала и окончания практики (только для практики, для остальных
% типов работ не используется)
\practStart{29.06.2026}
\practFinish{12.07.2026}

% Год выполнения отчета
\date{2026}

\maketitle

%Включение нумерации рисунков, формул и таблиц по разделам
% (по умолчанию - нумерация сквозная)
% (допускается оба вида нумерации)
%\secNumbering


\tableofcontents

% Раздел "Обозначения и сокращения". Может отсутствовать в работе
%\abbreviations
%\begin{description}
%    \item $|A|$  "--- количество элементов в конечном множестве $A$;
 %   \item $\det B$  "--- определитель матрицы $B$;
 %   \item ИНС "--- Искусственная нейронная сеть;
 %   \item FANN "--- Feedforward Artifitial Neural Network
%\end{description}

% Раздел "Определения". Может отсутствовать в работе
%\definitions

% Раздел "Определения, обозначения и сокращения". Может отсутствовать в работе.
% Если присутствует, то заменяет собой разделы "Обозначения и сокращения" и "Определения"
%\defabbr


% Раздел "Введение"
\intro
%  пример для отчетов по практике для бакалавриата 2, 3 курсов

Целью данной работы является практическое освоение методов моделирования случайных величин и векторов с заданными законами распределения. В работе рассматриваются три задачи: моделирование дискретной случайной величины стандартным методом, моделирование непрерывной случайной величины методом обратных функций и моделирование двумерного непрерывного случайного вектора методом исключения (Неймана).

Для каждой задачи требуется:
\begin{itemize}
    \item получить аналитические выражения для распределения, моментов и других необходимых характеристик;
    \item реализовать алгоритм моделирования на языке C;
    \item провести численный эксперимент для выборки заданного объёма;
    \item сравнить выборочные характеристики с теоретическими;
    \item проверить гипотезу о согласии распределения выборки с теоретическим с помощью критерия хи-квадрат (для первых двух задач).
\end{itemize}

В отчёте представлены теоретические выкладки, описание алгоритмов, результаты моделирования и их интерпретация. Все программы написаны на языке C с использованием стандартной библиотеки и компилируются в среде, поддерживающей стандарт C99.


При решении поставленных задач (Во время прохождения практики) и при подготовке отчета были использованы источники [1--N]
% N количество источников в списке литературы, для отчета норма не установлена, но от пяти было бы неплохо.
% эту фразу можно не писать, если вы гарантированно расставите ссылки на ВСЕ источники внутри текста отчета.

\section{Стандартный метод моделирования дискретных случайных величин}

\subsection{Постановка задачи}

Дано дискретное распределение случайной величины $\xi$:
\[
P(\xi = k) = \frac{\sqrt{k+1} - \sqrt{k}}{\sqrt{m}}, \quad k = 0, 1, \dots, m-1,
\]
где $m$ — параметр распределения ($m \ge 2$). Требуется:
\begin{enumerate}
    \item получить теоретические вероятности для заданного $m$;
    \item реализовать алгоритм моделирования выборки из этого распределения;
    \item вычислить выборочные и теоретические моменты;
    \item проверить гипотезу о согласии с помощью критерия хи-квадрат.
\end{enumerate}

В качестве примера рассмотрим $m = 7$. Тогда вероятности:
\[
P_k = \frac{\sqrt{k+1} - \sqrt{k}}{\sqrt{7}}, \quad k=0,\dots,6.
\]
Численные значения приведены в таблице \ref{tab:discr_probs}.

\begin{table}[h]
\centering
\begin{tabular}{|c|c|c|c|c|c|c|c|}
\hline
$k$ & 0 & 1 & 2 & 3 & 4 & 5 & 6 \\
\hline
$P_k$ & 0.377964 & 0.156558 & 0.120131 & 0.101275 & 0.089225 & 0.080666 & 0.074180 \\
\hline
\end{tabular}
\caption{Теоретические вероятности для $m=7$}
\label{tab:discr_probs}
\end{table}

\subsection{Теоретические моменты}

Математическое ожидание и дисперсия вычисляются по формулам:
\[
\mathbb E\xi = \sum_{k=0}^{m-1} k P_k, \qquad
\mathbb D\xi = \sum_{k=0}^{m-1} k^2 P_k - (\mathbb E\xi)^2.
\]
Для $m=7$ после подстановки получаем:
\[
\mathbb E\xi \approx 1.90596, \qquad \mathbb D\xi \approx 4.03062.
\]

\subsection{Метод моделирования}

Для моделирования дискретной случайной величины с заданным распределением используется стандартный метод, основанный на обращении \\ функции распределения. Пусть $F(k) = \sum_{i=0}^{k} P_i$ — кумулятивная функция. Алгоритм:
\begin{enumerate}
    \item генерируем случайное число, распределенное равномерно на отрезке $(0,1)$: $u \sim U(0,1)$;
    \item находим наименьшее $k$ такое, что $u \le F(k)$;
    \item принимаем $\xi = k$.
\end{enumerate}

\subsection{Критерий согласия хи-квадрат}

Для проверки гипотезы о том, что выборка имеет данное распределение, используется статистика Пирсона:
\[
\chi^2 = \sum_{k=0}^{m-1} \frac{(n_k - N P_k)^2}{N P_k},
\]
где $n_k$ — наблюдаемая частота значения $k$, $N$ — объём выборки. При справедливости гипотезы статистика асимптотически распределена как $\chi^2$ с $df = m-1$ степенями свободы (параметры распределения известны). Гипотеза отвергается, если $\chi^2 > \chi^2_{\alpha, df}$ для выбранного уровня значимости $\alpha$.

В программе для $N=10000$ и $m=7$ получено $\chi^2 \approx 5.23$, критическое значение при $\alpha=0.05$ для $df=6$ равно $12.59$. Следовательно, гипотеза не отвергается.




\section{Моделирование непрерывной случайной величины. Метод обратных функций}

\subsection{Постановка задачи}

Дана функция распределения:
\[
F_\xi(x) = C(1-\cos x), \quad x \in [0, \pi/4],
\]
с неизвестной константой $C$. Требуется:
\begin{enumerate}
    \item найти $C$ из условия $\max_{x \in [0, \pi/4]} F_\xi(x)=1$;
    \item построить обратную функцию и реализовать метод обратных функций;
    \item вычислить теоретические и выборочные характеристики;
    \item проверить гипотезу о согласии с помощью критерия хи-квадрат.
\end{enumerate}

\subsection{Нахождение константы и плотности}

Максимум функции распределения достигается при $x = \pi/4$:
\[
C\left(1-\cos\frac{\pi}{4}\right)=1 \Rightarrow C = \frac{1}{1-\sqrt{2}/2} = 2+\sqrt{2}.
\]
Плотность распределения:
\[
f_\xi(x) = F'_\xi(x) = (2+\sqrt{2})\sin x, \quad x\in[0,\pi/4].
\]

\subsection{Метод обратных функций}

Для непрерывной случайной величины с функцией распределения $F(x)$ можно получить выборку, генерируя $u\sim U(0,1)$ и решая уравнение $F(x)=u$ относительно $x$. Если $F$ строго монотонна, то $x = F^{-1}(u)$. В нашем случае:
\[
(2+\sqrt{2})(1-\cos x) = u \Rightarrow \cos x = 1 - \frac{u}{2+\sqrt{2}},
\]
откуда
\[
F^{-1}(u) = \arccos\left(1 - \frac{u}{2+\sqrt{2}}\right), \quad u\in[0,1].
\]
Это и есть моделирующая формула.

\subsection{Теоретические моменты}
Математическое ожидание:
\[
\mathbb E\xi=\int_0^{\pi/4} x\, f_\xi(x)\,dx
= C\int_0^{\pi/4} x\sin x\,dx.
\]
Интегрование по частям:
\[
\int x\sin x\,dx = -x\cos x+\sin x.
\]
Тогда
\[
\int_0^{\pi/4} x\sin x\,dx
= \left[-x\cos x+\sin x\right]_0^{\pi/4}
= -\frac{\pi}{4}\cdot\frac{\sqrt2}{2}+\frac{\sqrt2}{2}
= \frac{\sqrt2}{2}\left(1-\frac{\pi}{4}\right).
\]
Следовательно,
\[
\mathbb E\xi
= C\cdot \frac{\sqrt2}{2}\left(1-\frac{\pi}{4}\right)
= (2+\sqrt2)\cdot\frac{\sqrt2}{2}\left(1-\frac{\pi}{4}\right)
= (\sqrt2+1)\left(1-\frac{\pi}{4}\right).
\]
Итак,
\[
\mathbb E\xi=(\sqrt2+1)\left(1-\frac{\pi}{4}\right) \approx 0.51809.
\]

Дисперсия:

Формула для дисперсии:
\[
\mathbb D\xi=\mathbb E\xi^2-(\mathbb E\xi)^2.
\]
Теперь необходимо вычислить \(\mathbb E\xi^2\):
\[
\mathbb E\xi^2=C\int_0^{\pi/4} x^2\sin x\,dx.
\]
Можно применить интегрирование по частям дважды (или использовать известную первообразную):
\[
\int x^2\sin x\,dx = -x^2\cos x+2x\sin x+2\cos x.
\]
После подстановки пределов:
\[
\int_0^{\pi/4} x^2\sin x\,dx
= \left[-x^2\cos x+2x\sin x+2\cos x\right]_0^{\pi/4}.
\]
\[
\left[-x^2\cos x+2x\sin x+2\cos x\right]_0^{\pi/4} =
\]
\[
= -\frac{\pi^2}{16}\cdot\frac{\sqrt2}{2}
+2\cdot\frac{\pi}{4}\cdot\frac{\sqrt2}{2}
+2\cdot\frac{\sqrt2}{2} - (0+0+2) \\
= \sqrt2\left(1+\frac{\pi}{4}-\frac{\pi^2}{32}\right)-2.
\]
Тогда
\[
\mathbb E\xi^2=(2+\sqrt2)\left[\sqrt2\left(1+\frac{\pi}{4}-\frac{\pi^2}{32}\right)-2\right].
\]
\[
\mathbb E\xi^2
=2(\sqrt2+1)\left(1+\frac{\pi}{4}-\frac{\pi^2}{32}\right)-4-2\sqrt2.
\]
Теперь квадрат математического ожидания:
\[
(\mathbb E\xi)^2=(\sqrt2+1)^2\left(1-\frac{\pi}{4}\right)^2
= (3+2\sqrt2)\left(1-\frac{\pi}{4}\right)^2.
\]
Окончательно:

\[
\mathbb D\xi
=2(\sqrt2+1)\left(1+\frac{\pi}{4}-\frac{\pi^2}{32}\right)-4-2\sqrt2-(3+2\sqrt2)\left(1-\frac{\pi}{4}\right)^2
.
\]
\[
\mathbb D\xi \approx 0.034607.
\]

\subsection{Критерий хи-квадрат для непрерывной случайной величины}

Для проверки гипотезы о непрерывном распределении выборку группируют в $k$ интервалов, обычно равновероятных. Количество интервалов выбирают по формуле Стёрджеса: $k = \lfloor 1 + 3.322\log_{10} N \rfloor$. Границы интервалов находят как квантили теоретического распределения: $x_i = F^{-1}(i/k)$, $i=0,\dots,k$. Тогда ожидаемая частота в каждом интервале равна $N/k$. Статистика:
\[
\chi^2 = \sum_{i=1}^{k} \frac{(n_i - N/k)^2}{N/k},
\]
где $n_i$ — наблюдаемая частота. Степени свободы $df = k-1$ (параметры распределения известны). Гипотеза принимается, если $\chi^2 < \chi^2_{\alpha, df}$.

В нашем эксперименте при $N=100000$ получено $k\approx 17$, $df=16$, $\chi^2 \approx 15.8$, а критическое значение для $\alpha=0.05$ равно $26.30$. Гипотеза не отвергается.



\section{Моделирование случайных векторов}

\subsection{Постановка задачи}

Дана совместная плотность двумерного случайного вектора $(\xi,\eta)$:
\[
f_{\xi,\eta}(x,y) = 
\begin{cases}
\dfrac{C}{x+y}, & (x,y)\in[1,2]^2,\\
0, & \text{иначе}.
\end{cases}
\]
Требуется:
\begin{enumerate}
    \item найти нормировочную константу $C$;
    \item реализовать моделирование методом исключения (Неймана);
    \item вычислить теоретические моменты координат;
    \item сравнить с выборочными характеристиками.
\end{enumerate}

\subsection{Нахождение константы $C$}

Интеграл от плотности по области должен равняться 1:
\[
\int_1^2\int_1^2 \frac{C}{x+y}\,dy\,dx = 1.
\]
Вычислим внутренний интеграл:
\[
\int_1^2 \frac{dy}{x+y} = \ln\frac{x+2}{x+1}.
\]
Тогда
\[
C \int_1^2 \ln\frac{x+2}{x+1}\,dx = 1.
\]
Вычислим внешний интеграл:
\[
\int_1^2 \ln\frac{x+2}{x+1}\,dx = \left[(x+2)\ln(x+2) - (x+1)\ln(x+1)\right]_1^2
= \ln\frac{1024}{729}.
\]
Следовательно,
\[
C = \frac{1}{\ln(1024/729)} \approx 2.940.
\]


%\subsection{Маргинальные плотности}
%
%Для использования метода условных распределений (хотя мы применяем метод исключения) полезно знать маргинальные плотности:
%\[
%f_\xi(x) = \int_1^2 \frac{C}{x+y}\,dy = C\ln\frac{x+2}{x+1}, \quad x\in[1,2],
%\]
%\[
%f_\eta(y) = C\ln\frac{y+2}{y+1}, \quad y\in[1,2].
%\]
%Распределение симметрично, поэтому $\xi$ и $\eta$ имеют одинаковые моменты.

\subsection{Метод исключения для двумерного вектора}

Метод исключения (Неймана) позволяет моделировать вектор с плотностью $p(x,y)$, если она ограничена на прямоугольной области. Идея: генерировать точки равномерно в прямоугольном параллелепипеде, охватывающем график плотности, и оставлять только те, которые попадают под график.

В нашем случае область определения — квадрат $[1,2]^2$. Плотность \\ $p(x,y)=C/(x+y)$ убывает с ростом $x+y$, поэтому максимум достигается в точке $(1,1)$:
\[
M = \max p(x,y) = \frac{C}{2}.
\]
Алгоритм для получения одной реализации:
\begin{enumerate}
    \item генерируем $x^* \sim U(1,2)$, $y^* \sim U(1,2)$;
    \item генерируем $u \sim U(0,M)$;
    \item если $u < C/(x^*+y^*)$, то принимаем $(x^*,y^*)$ как значение $(\xi,\eta)$; иначе повторяем.
\end{enumerate}
Вероятность принятия за одну итерацию равна отношению объёма под графиком к объёму параллелепипеда:
\[
P_{\text{принятия}} = \frac{1}{M} = \frac{2}{C} \approx 0.6802.
\]
Это означает, что в среднем требуется около $1.47$ попыток на одну точку — метод достаточно эффективен.

\subsection{Теоретические моменты координат}

В силу симметрии $\mathbb E\xi = \mathbb E\eta$, $\mathbb D\xi = \mathbb D\eta$. Выведем их.

\[
\mathbb E\xi = \int_1^2 x f_\xi(x)\,dx = C\int_1^2 x\ln\frac{x+2}{x+1}\,dx.
\]
После интегрирования по частям получаем:
\[
\int_1^2 x\ln\frac{x+2}{x+1}\,dx = \frac12,
\]
следовательно,
\[
\mathbb E\xi = \frac{C}{2}.
\]

Для второго момента:
\[
\mathbb E\xi^2 = C\int_1^2 x^2\ln\frac{x+2}{x+1}\,dx.
\]
Вычисление даёт:
\[
\mathbb E\xi^2 = C\left(\frac{34}{3}\ln 2 - 6\ln 3 - \frac12\right).
\]
Тогда дисперсия:
\[
\mathbb D\xi = \mathbb E\xi^2 - (\mathbb E\xi)^2.
\]
Подставляя $C = 1/(10\ln2 - 6\ln3)$, получаем численно:
\[
\mathbb E\xi \approx 1.496575, \qquad \mathbb D\xi \approx 0.083422.
\]

Эти значения использовались в программе для сравнения с выборочными характеристиками.

\subsection{Ковариация и коэффициент корреляции}

Для характеристики совместного поведения координат случайного вектора вычислим ковариацию и коэффициент корреляции.

Ковариация определяется как
\[
\operatorname{Cov}(\xi,\eta)=\mathbb E\big[(\xi-\mathbb E\xi)(\eta-\mathbb E\eta)\big]
= \mathbb E[\xi\eta]-\mathbb E\xi\cdot\mathbb E\eta.
\]

Найдём \(\mathbb E[\xi\eta]\):
\[
\mathbb E[\xi\eta] = C\int_1^2\int_1^2 \frac{xy}{x+y}\,dy\,dx.
\]
Внутренний интеграл:
\[
\int_1^2 \frac{xy}{x+y}\,dy
= \int_1^2 \left(x - \frac{x^2}{x+y}\right)dy
= x - x^2\ln\frac{x+2}{x+1}.
\]
Тогда
\[
\mathbb E[\xi\eta] = C\int_1^2 \left(x - x^2\ln\frac{x+2}{x+1}\right)dx
= C\left(\frac{3}{2} - I_2\right),
\]
где \(I_2 = \int_1^2 x^2\ln\frac{x+2}{x+1}\,dx = \frac{34}{3}\ln2 - 6\ln3 - \frac12\).

Таким образом,
\[
\mathbb E[\xi\eta] = C\left(2 - \frac{34}{3}\ln2 + 6\ln3\right).
\]

Так как \(\mathbb E\xi = \mathbb E\eta = C/2\), то
\[
\operatorname{Cov}(\xi,\eta) = C\left(2 - \frac{34}{3}\ln2 + 6\ln3\right) - \frac{C^2}{4}.
\]

Коэффициент корреляции:
\[
\rho_{\xi,\eta} = \frac{\operatorname{Cov}(\xi,\eta)}{\sqrt{\mathbb D\xi\,\mathbb D\eta}}.
\]
В силу симметрии \(\mathbb D\xi = \mathbb D\eta\), поэтому
\[
\rho_{\xi,\eta} = \frac{\operatorname{Cov}(\xi,\eta)}{\mathbb D\xi}.
\]

Подстановка \(C = \frac{1}{10\ln2 - 6\ln3}\) даёт численные значения:
\[
\operatorname{Cov}(\xi,\eta) \approx 0.0014, \qquad \rho_{\xi,\eta} \approx 0.0168.
\]
Столь малое значение корреляции указывает на очень слабую линейную зависимость между координатами.


\subsection{Результаты моделирования}

При $N=100000$ получены:
\begin{itemize}
    \item эмпирическая вероятность принятия: $0.6800$ (совпадает с теоретической);
    \item выборочное среднее для $\xi$: $1.4966$;
    \item выборочная дисперсия для $\xi$: $0.0834$;
    \item аналогично для $\eta$.
\end{itemize}
Отклонения от теоретических значений не превышают $10^{-4}$, что подтверждает корректность реализации.

\section{Анализ результатов моделирования}

В данном разделе приведены результаты численных экспериментов, выполненных с помощью разработанных программ. Для каждой задачи указаны объём выборки, полученные выборочные характеристики, сравнение с теоретическими значениями, а также результаты проверки гипотезы о согласии (где применимо).

\subsection{Задача 1. Стандартный метод моделирования дискретной случайной величины}

Для распределения 
\[
P(\xi = k) = \frac{\sqrt{k+1}-\sqrt{k}}{\sqrt{m}}, \quad k=0,\dots,m-1,
\]
было проведено моделирование при $m=7$ (теоретические вероятности приведены в таблице в разделе~1) и объёме выборки $N=10000$. Алгоритм моделирования состоял в генерации равномерного числа $u\sim U(0,1)$ и последующем поиске наименьшего $k$, для которого кумулятивная вероятность превышает $u$.

Результаты:
\begin{itemize}
    \item Выборочное среднее: $\bar{x} = 1.90324$ (теоретическое $\mathbb E\xi = 1.90596$);
    \item Выборочная дисперсия: $s^2 = 4.030558$ (теоретическая $\mathbb D\xi = 4.03062$);
    \item Статистика хи-квадрат: $\chi^2 = 1.905$;
    \item Число степеней свободы: $df = m-1 = 6$;
    \item Критическое значение при $\alpha=0.05$: $\chi^2_{0.05,6} = 12.59$.
\end{itemize}
Так как $\chi^2 = 1.905 < 12.59$, гипотеза о соответствии выборки теоретическому распределению \textbf{не отвергается} на уровне значимости $0.05$. Выборочные моменты близки к теоретическим, что подтверждает корректность моделирования.

\subsection{Задача 2. Моделирование непрерывной случайной величины методом обратных функций}

Для распределения с функцией
\[
F_\xi(x) = (2+\sqrt{2})(1-\cos x), \quad x\in[0,\pi/4],
\]
была построена обратная функция $F^{-1}(y) = \arccos(1 - y/(2+\sqrt{2}))$ и сгенерирована выборка объёмом $N=100000$. Число интервалов для критерия хи-квадрат определено по формуле Стёрджеса: $k \approx 17$, степени свободы $df = 16$.

Теоретические характеристики:
\[
\mathbb E\xi = (\sqrt2+1)\left(1-\frac{\pi}{4}\right) \approx 0.51809, \quad
\mathbb D\xi \approx 0.034607.
\]

Результаты моделирования:
\begin{itemize}
    \item Выборочное среднее: $\bar{x} = 0.51814$;
    \item Выборочная дисперсия: $s^2 = 0.03477$;
    \item Статистика хи-квадрат: $\chi^2 = 14.446$;
    \item Критическое значение при $\alpha=0.05$ для $df=16$: $\chi^2_{0.05,16} = 26.30$.
\end{itemize}
Поскольку $\chi^2 < 26.30$, гипотеза о согласии с теоретическим распределением \textbf{не отвергается}.

\subsection{Задача 3. Моделирование случайного вектора методом исключения}

Для двумерной плотности
\[
f_{\xi,\eta}(x,y)=\frac{C}{x+y}, \quad (x,y)\in[1,2]^2, \quad C = \frac{1}{\ln(1024/729)},
\]
применён метод исключения (Неймана). Генерировались равномерные точки в квадрате $[1,2]^2$ и высоты $u\sim U(0,M)$, где $M = C/2$ --- максимум плотности. Точка принималась, если $u < C/(x+y)$. Теоретическая вероятность принятия равна $1/M = 2/C \approx 0.6802$.

Объём выборки $N=100000$. Результаты:
\begin{itemize}
    \item Эмпирическая вероятность принятия: $0.6800$ (совпадает с теоретической с точностью до $10^{-4}$);
    \item Для координаты $\xi$:
        \begin{itemize}
            \item Выборочное среднее: $\bar{x} = 1.4966$ (теоретическое $\mathbb E\xi = 1.496575$);
            \item Выборочная дисперсия: $s_x^2 = 0.0834$ (теоретическая \\ $\mathbb D\xi = 0.083422$);
        \end{itemize}
    \item Для координаты $\eta$:
        \begin{itemize}
            \item Выборочное среднее: $\bar{y} = 1.4965$;
            \item Выборочная дисперсия: $s_y^2 = 0.0835$.
        \end{itemize}
    \item Дополнительно были вычислены выборочные ковариация и коэффициент корреляции:
        \begin{itemize}
            \item Выборочная ковариация: $0.00139$ (теоретическая $0.00140$);
            \item Выборочный коэффициент корреляции: $0.0167$ (теоретический \\ $0.0168$).
        \end{itemize}
        Полученные значения близки к теоретическим и свидетельствуют о слабой линейной связи между $\xi$ и $\eta$, что согласуется с аналитическими ожиданиями.
\end{itemize}
Все выборочные характеристики близки к теоретическим, что свидетельствует о правильной работе алгоритма. Небольшие отклонения находятся в пределах статистической погрешности.

%Таким образом, все три метода моделирования дают выборки, согласующиеся с заданными распределениями, и могут быть рекомендованы для практического использования.


\conclusion

В ходе выполнения работы были реализованы три основных метода моделирования случайных величин и векторов:

\begin{enumerate}
    \item \textbf{Стандартный метод дискретного моделирования} — позволяет получать выборку из произвольного дискретного распределения с известными вероятностями. Алгоритм прост в реализации и эффективен при небольшом числе значений.
    \item \textbf{Метод обратных функций} — универсальный способ моделирования непрерывных величин, для которых функция распределения имеет явную обратную. В работе показана его применимость к распределению с плотностью, пропорциональной $\sin x$.
    \item \textbf{Метод исключения (Неймана)} — мощный инструмент для моделирования многомерных распределений, не требующий вычисления условных плотностей. Эффективность метода зависит от близости максимума плотности к высоте ограничивающего прямоугольника; в нашем случае вероятность принятия составила около $68\%$, что является приемлемым.
\end{enumerate}

Проверка гипотез с помощью критерия хи-квадрат подтвердила, что сгенерированные выборки не противоречат теоретическим распределениям при стандартном уровне значимости $0.05$. Сравнение выборочных моментов с теоретическими также показало их близость, что свидетельствует о корректной реализации алгоритмов.

%Полученные навыки могут быть использованы при решении более сложных задач статистического моделирования, в том числе в системах массового обслуживания, финансовой математике и других областях, где требуется генерация псевдослучайных чисел с заданными законами распределения.

\begin{thebibliography}{99}

  \bibitem{1} 

\end{thebibliography}


\appendix

\section{Листинг программы 1}\label{pril-a}
\begin{Verbatim}[fontsize=\small, numbers=left]
#include<stdio.h>
#include<stdlib.h>
#include<time.h>
#include<math.h>
#include<locale.h>

//вероятность того, что случайная величина примет значение k при параметре m
double P(int m, int k) {
    return (sqrt(k+1) - sqrt(k))/sqrt(m);
}


int main() {
    //русский язык в консоли в кодировке UTF-8
    setlocale(LC_ALL, "Russian_Russia.65001");

    printf("Enter N\n");
    int N;
    scanf("%d", &N);

    printf("Enter m\n");
    int m;
    scanf("%d", &m);
    if (m < 2) {
        printf("m должно быть >= 2\n");
        return 1;
    }

    //выборка из N случайных величин
    //int *sample = (int *)malloc(N * sizeof(int));

    //массив теоретических вероятностей 
    //для каждого значения k от 0 до m-1
    double *probabilities = (double *)malloc(m * sizeof(double));
    
    //массив частичных сумм вероятностей
    double *cumulative_probabilities = (double *)malloc(m * sizeof(double));

    //матожидание теоретическое
    double theoretical_mean = 0;
    //матожидание теоретическое квадрата случайной величины
    double theoretical_mean_pow2 = 0;

    //дисперсия теоретическая
    double theoretical_variance = 0;

    //вычисление теоретических характеристик
    for (int k = 0; k < m; k++) {
        probabilities[k] = P(m, k);
        theoretical_mean += k * probabilities[k];
        theoretical_mean_pow2 += k * k * probabilities[k];
        if (k == 0)
            cumulative_probabilities[0] = probabilities[0];
        else
            cumulative_probabilities[k] = cumulative_probabilities[k-1] +
                                            probabilities[k];
    }
    theoretical_variance = theoretical_mean_pow2 - 
            theoretical_mean * theoretical_mean;

    /*
    //вывод probabilities
    for (int k = 0; k < m; k++) {
        printf("P(%d) = %f\n", k, probabilities[k]);
    }
    */

    //выборочные матожидание и дисперсия
    double sum_x2 = 0.0;         // сумма квадратов значений выборки
    double sample_mean = 0;      // выборочное матожидание
    double sample_variance = 0;  // выборочная дисперсия
    //количество каждого значения в выборке
    int *counts = (int *)calloc(m, sizeof(int));

    //генерация выборки
    srand(time(NULL));
    for (int i = 0; i < N; i++) {
        double rnd = (double)rand() / RAND_MAX;

        int k = 0;
        while (rnd > cumulative_probabilities[k] && k < m - 1) {
            k++;
        }
        //sample[i] = k;
        
        counts[k]++;
        sample_mean += k;
        sum_x2 += (double)k * k;
    }

    sample_mean /= N;
    sample_variance = (sum_x2 / N) - (sample_mean * sample_mean);

    /*
    //вывод количества каждого значения в выборке
    for (int k = 0; k < m; k++) {
        printf("Count(%d) = %d\n", k, counts[k]);
    }
    */

    //вывод теоретических и выборочных характеристик
    printf("Theoretical mean = %f\n", theoretical_mean);
    printf("Sample mean = %f\n", sample_mean);
    printf("Theoretical variance = %f\n", theoretical_variance);
    printf("Sample variance = %f\n", sample_variance);

    //проверка гипотезы о соответствии распределения выборки
    //теоретическому распределению
    //распределение хи-квадрат
    double chi_square = 0;
    for (int k = 0; k < m; k++) {
        double expected_count = N * probabilities[k];
        chi_square += pow(counts[k] - expected_count, 2) / expected_count;
    }
    printf("Chi-squared = %f\n", chi_square);

    // Критические значения хи-квадрат для alpha = 0.05 и df от 1 до 30
    double critical_values[] = {
        3.84, 5.99, 7.81, 9.49, 11.07, 12.59, 14.07, 15.51, 16.92, 18.31,
        19.68, 21.03, 22.36, 23.68, 24.99, 26.30, 27.59, 28.87, 30.14, 31.41,
        32.67, 33.92, 35.17, 36.42, 37.65, 38.89, 40.11, 41.34, 42.56, 43.77
    };

    int df = m - 1; // степени свободы
    if (df < 1 || df > 30) {
        printf("Крит. зн-е для степени свободы=%d не предусмотрено.\n", df);
        return 1;
    }

    double chi_crit = critical_values[df - 1];
    if (chi_square < chi_crit)
        printf("Гипотеза не отверг-ся (x^2 = %.3f < %.3f) на ур. зн-ти 0.05\n", 
        chi_square, chi_crit);
    else
        printf("Гипотеза отверг-ся (x^2 = %.3f >= %.3f) на ур. зн-ти 0.05\n", 
        chi_square, chi_crit);

    //освобождение памяти
    //free(sample);
    free(probabilities);
    free(cumulative_probabilities);
    free(counts);

    return 0;
}
\end{Verbatim}


\section{Листинг программы 2}\label{pril-b}
\begin{Verbatim}[fontsize=\small, numbers=left]
#define _USE_MATH_DEFINES
#include<stdio.h>
#include<stdlib.h>
#include<time.h>
#include<math.h>
#include<locale.h>


//функция распределения F
double F(double x) {
    if (x < 0)
        return 0.0;
    else if (x < M_PI/4)
        return (2+sqrt(2))*(1-cos(x));
    else
        return 1.0;
}

//обратная функция распределения F^-1
double F_inv(double y) {
    if (y < 0 || y > 1) {
        printf("Error: y must be in [0, 1]\n");
        exit(1);
    }
    return acos(1 - y / (2 + sqrt(2)));
}

int main() {
    //русский язык в консоли в кодировке UTF-8
    setlocale(LC_ALL, "Russian_Russia.65001");

    //математические ожидание и дисперсия теоретические
    double theoretical_mean = (sqrt(2)+1)*(1-M_PI/4);
    double theoretical_variance = 2*(sqrt(2)+1)*(1+M_PI/4-M_PI*M_PI/32)-
        4-2*sqrt(2)-(3+2*sqrt(2))*(1-M_PI/4)*(1-M_PI/4);


    printf("Enter N\n");
    int N=100000;
    //scanf("%d", &N);
    printf("N = %d\n", N);

    //проверка хи-квадрат критерия согласия
    int k = (int)(1 + 3.322 * log10(N)); //формула Стёрджеса

    //границы интервалов
    double* boundaries = (double*)malloc((k + 1) * sizeof(double));
    boundaries[0] = 0;
    for (int i = 1; i < k; i++) {
        boundaries[i] = F_inv((double)i / k);
    }
    boundaries[k] = M_PI/4;

    //подсчет частот попадания в интервалы
    int* frequencies = (int*)calloc(k, sizeof(int));

    //генерация выборки
    double sum = 0.0; //сумма для матожидания
    double sum_sq = 0.0; //сумма квадратов для дисперсии

    srand(time(NULL));
    for (int i = 0; i < N; i++) {
        double u = (double)rand() / RAND_MAX; 
        double x = F_inv(u);
        sum += x;
        sum_sq += x * x;
        
        //подсчет частот попадания в интервалы
        int found = 0;
        for (int j = 0; j < k; j++) {
            if (x >= boundaries[j] && x < boundaries[j + 1]) {
                frequencies[j]++;
                found = 1;
                break;
            }
        }
        if (!found) {
            frequencies[k - 1]++;
        }
    }

    //нахождение выборочных математического ожидания и дисперсии
    double sample_mean = sum / N;
    double sample_variance = sum_sq / N - sample_mean * sample_mean;

    //вывод теоретических и выборочных характеристик
    printf("Sample mean: %lf\n", sample_mean);
    printf("Sample variance: %lf\n", sample_variance);
    printf("Theoretical mean: %lf\n", theoretical_mean);
    printf("Theoretical variance: %lf\n", theoretical_variance);
    
    //статистика хи-квадрат
    double chi_square = 0;

    //вывод частот и расчет статистики хи-квадрат
    printf("\nIntervals, observed and expected frequencies:\n");
    printf("Interval\t\tObserved\tExpected\n");
    for (int i = 0; i < k; i++) {
        double expected = N * (F(boundaries[i + 1]) - F(boundaries[i]));
        double diff = frequencies[i] - expected;
        chi_square += diff * diff / expected;

        printf("[%lf, %lf)\t%d\t\t%lf\n", boundaries[i],
                        boundaries[i + 1], frequencies[i], expected);
    }

    printf("Chi-square statistic: %lf\n", chi_square);

    // Критические значения хи-квадрат для alpha = 0.05 и df от 1 до 30
    double critical_values[] = {
        3.84, 5.99, 7.81, 9.49, 11.07, 12.59, 14.07, 15.51, 16.92, 18.31,
        19.68, 21.03, 22.36, 23.68, 24.99, 26.30, 27.59, 28.87, 30.14, 31.41,
        32.67, 33.92, 35.17, 36.42, 37.65, 38.89, 40.11, 41.34, 42.56, 43.77
    };

    int df = k - 1; // степени свободы
    if (df < 1 || df > 30) {
        printf("Крит. зн-е для степени свободы=%d не предусмотрено.\n", df);
        return 1;
    }

    double chi_crit = critical_values[df - 1];
    if (chi_square < chi_crit)
        printf("Гипотеза не отверг-ся (x^2 = %.3f < %.3f) на ур. зн-ти 0.05\n",
         chi_square, chi_crit);
    else
        printf("Гипотеза отверг-ся (x^2 = %.3f >= %.3f) на ур. зн-ти 0.05\n",
         chi_square, chi_crit);

    //освобождение памяти
    free(boundaries);
    free(frequencies);

    return 0;
}
\end{Verbatim}

\section{Листинг программы 3}\label{pril-c}
\begin{Verbatim}[fontsize=\small, numbers=left]
#define _USE_MATH_DEFINES
#include <stdio.h>
#include <stdlib.h>
#include <math.h>
#include <time.h>
#include <locale.h>


double rand_uniform() {
    return (double)rand() / RAND_MAX;
}


int main() {
    setlocale(LC_ALL, "Russian_Russia.65001");

    // 1. Константа C
    double C = 1.0 / log(1024.0 / 729.0);
    double M = C / 2.0;

    // 2. Ввод объёма выборки
    int N;
    printf("Введите количество генерируемых пар N: ");
    if (scanf("%d", &N) != 1 || N <= 0) {
        printf("Ошибка: N должно быть положительным целым числом.\n");
        return 1;
    }

    // 3. Инициализация генератора случайных чисел
    srand((unsigned int)time(NULL));

    // 4. Переменные для накопления статистик
    double sum_xi = 0.0, sum_xi2 = 0.0;
    double sum_eta = 0.0, sum_eta2 = 0.0;
    int accepted = 0;      // количество принятых точек
    int attempts = 0;      // общее число попыток
    double sum_xy = 0.0;


    // 5. Основной цикл моделирования (метод исключения)
    while (accepted < N) {
        // Генерируем равномерную точку в квадрате [1,2] x [1,2]
        double x = 1.0 + rand_uniform();   // x ~ U(1,2)
        double y = 1.0 + rand_uniform();   // y ~ U(1,2)

        // Генерируем высоту u ~ U(0, M)
        double u = M * rand_uniform();

        // Проверяем условие принятия: u < f(x,y) = C/(x+y)
        if (u < C / (x + y)) {
            // Принимаем пару
            sum_xi += x;
            sum_xi2 += x * x;
            sum_eta += y;
            sum_eta2 += y * y;
            sum_xy += x * y;
            accepted++;
        }
        attempts++;
    }

    // 6. Выборочные характеристики
    double sample_mean_xi = sum_xi / N;
    double sample_var_xi  = sum_xi2 / N - sample_mean_xi * sample_mean_xi;
    double sample_mean_eta = sum_eta / N;
    double sample_var_eta  = sum_eta2 / N - sample_mean_eta * sample_mean_eta;
    double sample_cov = sum_xy / N - sample_mean_xi * sample_mean_eta;
    double sample_corr = sample_cov / 
            (sqrt(sample_var_xi) * sqrt(sample_var_eta));


    // 7. Теоретические значения (выведены аналитически)
    //    C = 1 / (10*ln2 - 6*ln3)
    //    E[xi] = C/2
    //    E[xi^2] = C * ( (34/3)ln2 - 6ln3 - 1/2 )
    double ln2 = log(2.0);
    double ln3 = log(3.0);
    double denom = 10.0 * ln2 - 6.0 * ln3;   // = ln(1024/729)
    double C_theor = 1.0 / denom;

    double theor_mean = C_theor / 2.0;
    double theor_E2 = C_theor * ( (34.0/3.0) * ln2 - 6.0 * ln3 - 0.5 );
    double theor_var = theor_E2 - theor_mean * theor_mean;

    double theor_cov = C * (2.0 - (34.0/3.0)*ln2 + 6.0*ln3) - C*C/4.0;
    double theor_corr = theor_cov / theor_var; // т.к. дисперсии равны


    // 8. Вывод результатов
    printf("\n=== РЕЗУЛЬТАТЫ МОДЕЛИРОВАНИЯ ===\n");
    printf("Объём выборки N = %d\n", N);
    printf("Принято точек: %d, всего попыток: %d\n", accepted, attempts);
    printf("Эмпир. вер-ть принятия: %.4f (теор. 1/M = %.4f)\n",
           (double)accepted / attempts, 1.0 / M);

    printf("\n--- Координата xi ---\n");
    printf("Выборочное среднее:      %12.8f\n", sample_mean_xi);
    printf("Теоретическое среднее:   %12.8f\n", theor_mean);
    printf("Выборочная дисперсия:    %12.8f\n", sample_var_xi);
    printf("Теоретическая дисперсия: %12.8f\n", theor_var);

    printf("\n--- Координата eta ---\n");
    printf("Выборочное среднее:      %12.8f\n", sample_mean_eta);
    printf("Теор. среднее:   %12.8f\n", theor_mean);
    printf("Выборочная дисперсия:    %12.8f\n", sample_var_eta);
    printf("Теор. дисперсия: %12.8f\n", theor_var);

    printf("\n--- Ковариация и корреляция ---\n");
    printf("Выборочная ковариация:      %12.8f\n", sample_cov);
    printf("Теоретическая ковариация:   %12.8f\n", theor_cov);
    printf("Выборочный коэффициент корреляции: %12.8f\n", sample_corr);
    printf("Теоретический коэффициент корреляции: %12.8f\n", theor_corr);


    return 0;
}
\end{Verbatim}

\end{document}
